% Generated draft: requires mathematical and source review.
\documentclass[11pt]{article}
\usepackage[margin=25mm]{geometry}
\usepackage{fontspec}
\setmainfont{Latin Modern Roman}
\usepackage{amsmath,amssymb,mathtools,unicode-math}
\setmathfont{Latin Modern Math}
\usepackage{enumitem,needspace,xurl,hyperref}
\setlength{\parindent}{0pt}
\setlength{\parskip}{5pt}
\setlength{\emergencystretch}{3em}
\newcommand{\N}{\mathbb{N}}
\newcommand{\Z}{\mathbb{Z}}
\newcommand{\Q}{\mathbb{Q}}
\newcommand{\R}{\mathbb{R}}
\newcommand{\C}{\mathbb{C}}
\newcommand{\F}{\mathbb{F}}
\newcommand{\A}{\mathbb{A}}
\newcommand{\PP}{\mathbb{P}}
\newcommand{\E}{\mathbb{E}}
\newcommand{\eps}{\varepsilon}
\newcommand{\dd}{\,\mathrm{d}}
\begin{document}
\section*{UnsolvedMath ID: 11100023}
\textbf{Source problem code:} catalogue entry
\subsection{Definitions and mathematical statement}
\emph{Elliptic-cohomology research question (Mark Hovey's list).} Can the proposal of Dennis McLaughlin and Jean-Luc Brylinski to use gerbes, or perhaps \(2\)-gerbes, be developed into the intended geometric approach to elliptic cohomology? The antecedent ``along these lines'' is omitted, and the source expressly does not choose between gerbes and \(2\)-gerbes or specify the elliptic cohomology theory.

A gerbe is conventionally described by a suitable stack of groupoids that is locally nonempty and locally has isomorphic objects, with its band, topology, and additional geometry specified. This is more precise than the source author's tentative phrase ``a bundle of groupoids.'' A \(2\)-gerbe requires corresponding higher categorical data. The excerpt fixes none of those conventions. A precise proposed construction must specify the geometric objects, equivalences, pullbacks, degree and coefficient data, and how they represent the intended elliptic-cohomology classes or structures. The existence of gerbes alone does not furnish such a theory, and no particular modern spectrum is silently selected as the target.

The source's comment that Hovey had heard nothing for a while is historical. The supplied research summary describes gerbes and \(2\)-gerbes as relevant to string structures, \(\sigma\)-orientations and \(TMF\), mentions Stolz--Teichner and derived/categorical approaches and Brylinski--McLaughlin cocycles, and calls the literal proposal partly superseded. It also says no completed construction directly from gerbe cocycles was found, but its discussion is truncated. Those contextual successors do not establish the original construction. Its claims about a transgressed degree-\(2\) class and realization of the broader goal omit the class, transgression convention, hypotheses, and precise theorem, and are not adopted as verified statements.
\subsection{Short English statement}
Develop the McLaughlin--Brylinski gerbe or \(2\)-gerbe idea into a precise geometric elliptic-cohomology construction. The source leaves the target theory and geometric data open; modern string/TMF-related applications are context, not a supplied realization of this proposal.
\subsection{Sources}
\begin{itemize}
\item Local UnsolvedMath registry, permanent ID 11100023, source code catalogue entry; proposer Mark Hovey, ``Elliptic cohomology,'' item 3. The source explicitly names Dennis McLaughlin and Jean-Luc Brylinski.
\item \emph{Hovey's Algebraic Topology Problem List}, surviving mirror: \url{https://www-users.cse.umn.edu/~tlawson/hovey/}. Recorded access 2026-07-29; extraction html-ordered-list. The mirror has 65 top-level items across nine category pages.
\item The supplied summary mentions Brylinski, Brylinski--McLaughlin, string structures, \(\sigma\)-orientations to \(TMF\), Stolz--Teichner, and derived/categorical descriptions, without full citations. Its final no-completed-gerbe-cocycle-construction discussion is truncated.
\item Supplied status partially\_solved and PARTIAL-PROGRESS. The summary says partial/largely superseded with real successors, but provides no exact construction answering the literal request. Current primary-literature status and redistribution terms remain NEEDS\_REVIEW.
\end{itemize}
\end{document}
